% §17.6 Maximizing Confidentiality
%
% Monograph register: confidentiality as the 1-PPV posterior; the ch7 PPV
% formula IS the confidentiality formula; trading rate-budget for
% confidentiality on the perfect-filter line; worked example. Tightened.
% Folds salvaged ex-11 as rem:confidentiality-bits (information-theoretic
% leakage L = H(pi) - E[H(posterior)] in bits/query, complementing
% C = 1 - PPV). Cross-refs def:perfect-filter-line, def:ppv, def:npv,
% def:cipher-map, sec:trapdoor-paradigm, sec:fhe-sidebar. DEFINED:
% sec:maximizing-confidentiality, rem:confidentiality-bits.

\section{Maximizing Confidentiality}
\label{sec:maximizing-confidentiality}

\Cref{sec:trapdoor-paradigm} named a duality the framework has not yet
quantified: the cipher-map rates $\fprate$ and $\fnrate$ govern both
correctness and confidentiality. Looser rates make the server's output
less faithful to $f$ but also less informative to an adversary inspecting
it. The earlier sections treated correctness; this one treats
confidentiality, and the question is concrete: given a cipher map at fixed
$(\fprate, \fnrate)$ on a workload of known prevalence $\pi$, how much
does the output reveal about the plaintext? The chapter 7 PPV machinery
answers it with no further apparatus.

\subsection*{Confidentiality is a posterior}

An adversary watching the server's outputs faces a posterior-inference
problem: each output is a noisy report about the plaintext, and Bayes'
rule combines it with prior knowledge. For a single Boolean query the
target is
\[
  \Pr\bigl[\,x \in A \;\big\vert\; \tilde f(x) = \top\,\bigr],
\]
the probability that a reported $\top$ identifies a genuine member rather
than a false positive. Close to $1$, the report has pinned membership
down; close to the prior $\pi$, it has revealed nothing. The
confidentiality is the gap, which we take to be
\[
  C \;=\; 1 - \PPV,
\]
the share of $\top$-reports that are false positives. Confidentiality is
high when most $\top$-reports do not identify members and zero in the FHE
limit, where every report is correct and membership is fully exposed.

The posterior is exactly the positive predictive value of
\Cref{def:ppv}. Reading the chapter 7 formula in the cipher-map setting,
\[
  \PPV = \frac{\pi(1 - \fnrate)}{\pi(1 - \fnrate) + (1 - \pi)\fprate},
  \qquad
  C = 1 - \PPV = \frac{(1 - \pi)\fprate}{\pi(1 - \fnrate) + (1 - \pi)\fprate},
\]
the numerator $\pi(1-\fnrate)$ the rate of correct $\top$ on a member,
$(1-\pi)\fprate$ the rate of spurious $\top$ on a non-member, $C$ the
share that protects the plaintext. Zero false positives gives $C = 0$; an
unconditional $\top$ gives $C = 1 - \pi$. The negative case mirrors it,
the confidentiality of a $\bot$-report being $1 - \NPV$ (\Cref{def:npv}),
though in typical encrypted search $\pi$ is small and the $\top$-side
ambiguity dominates.

The $C = 1 - \PPV$ measure is the per-report complement; the
information-theoretic complement, in bits, is also available.

\begin{remark}[Confidentiality in bits per query]
\label{rem:confidentiality-bits}
Model the adversary's per-query posterior as Bernoulli at $\PPV$ for a
$\top$-report and at $1 - \NPV$ for a $\bot$-report. The leakage is the
mutual information between the plaintext bit and the report,
\[
  L \;=\; H(\pi) - \mathbb{E}_{\text{report}}\!\left[H(\text{posterior})\right]
  \;=\; H(\pi) - \bigl[\Pr(\top)\,H(\PPV) + \Pr(\bot)\,H(1 - \NPV)\bigr],
\]
with $H$ the binary entropy and $\Pr(\top) = \pi(1-\fnrate) +
(1-\pi)\fprate$. In the FHE limit $\fprate = \fnrate = 0$ the posteriors
are deterministic, $L = H(\pi)$, the full prior entropy leaks. A cipher
map with positive rates leaks a fraction of $H(\pi)$ per query, the
information-theoretic counterpart of the $C = 1 - \PPV$ posterior, and the
chapter 8 entropy machinery measures that fraction directly.
\end{remark}

\subsection*{Trading rate-budget for confidentiality}

On the perfect-filter line of \Cref{def:perfect-filter-line} the designer
faces a one-dimensional trade-off. An $\fnrate = 0$ map sits on the
$\fprate$-axis, and raising $\fprate$ raises $(1-\pi)\fprate$ and so
raises $C$ monotonically. Walking the line away from the origin trades
correctness for confidentiality. The trade is not free in the convenient
direction: chapter 5's space formula makes storage a decreasing function
of $\fprate$, so \emph{lower-storage cipher maps are more confidential},
one of the framework's small surprises, following entirely from the
chapter 7 posterior with no cryptographic argument. Storage couples to
accuracy through the chapter 8 entropy bound and accuracy to
confidentiality through the chapter 7 posterior, so one parameter sets
three coupled costs. Maximizing $C$ in the abstract is degenerate (walk
to $\fprate = 1$, $C = 1 - \pi$, the map useless), so the real problem
constrains the operational error: maximize $C$ subject to $\fprate \le
\fprate^*$ for an application-imposed cap, the optimum being the cap
itself.

\subsection*{A worked example}

Take prevalence $\pi = 0.01$ and rates $\fprate = 0.05$, $\fnrate = 0$.
The PPV formula gives
\[
  \PPV = \frac{0.01}{0.01 + 0.99 \cdot 0.05} = \frac{0.01}{0.0595}
  \approx 0.168,
  \qquad C \approx 0.832.
\]
When the server reports $x \in A$, the adversary's posterior assigns about
$16.8\%$ to genuine membership and $83.2\%$ to a false positive; one query
leaves the question far from settled. The FHE baseline reports with zero
error, $\PPV = 1$, $C = 0$: every report is correct and fully revealing.
The contrast is stark, the cipher map paying a $5\%$ false-positive rate
for an $83\%$ confidentiality margin, FHE paying a $10^4\times$ space and
$10^9\times$ time overhead for a zero margin. The example also shows the
workload dependence: holding $(\fprate, \fnrate) = (0.05, 0)$ but moving
to a balanced workload $\pi = 0.5$ gives $\PPV = 0.5/0.525 \approx 0.952$
and $C \approx 0.048$. The same map is confidential on a rare-positive
workload and revealing on a balanced one. The framework keeps the
parameters separate: the map carries $(\fprate, \fnrate)$, the workload
carries $\pi$, and $C(\fprate, \fnrate, \pi)$ is one function the designer
plots and reads. Chapter 7 did the algebra; chapter 17 inherits it.
